% GENERATED FILE. Edit canonical lessons, then run npm run generate:books.
% canonical-source-hash: fnv1a64:b309b00e0edaace8
% canonical-lessons: GU-C270-svabhav, GU-C270-dayalu, GU-C270-pramanik, GU-C270-ramuji, GU-C270-ghamandi

\chapter{Kind, honest and funny}
\label{ch:gu-kind-honest-and-funny}
\hlchaptermodality{270}

\begin{chapteropening}
\textbf{By the end of this chapter:} \emph{I can say what people are like and what they are not like.}

Complete the last lesson of chapter 270: I can say what people are like and what they are not like.
\end{chapteropening}

\section[\texorpdfstring{\gu{સ્વભાવ} (svabhāv)}{સ્વભાવ (svabhāv)}]{\gu{સ્વભાવ} (svabhāv) — temperament, nature}

\label{lesson:GU-C270-svabhav}



\begin{quote}
Before the new one: say the Gujarati for to remember fondly, then the Gujarati for respect.
\end{quote}

\subsection*{You'll want to know: \gu{સ્વભાવ}}
\textbf{\gu{સ્વભાવ}} — \emph{svabhāv} — \textquotedblleft{}temperament, nature\textquotedblright{}.

\begin{quote}
\textbf{\gu{તેનો} \gu{સ્વભાવ} \gu{શાંત} \gu{નથી}.} — \emph{teno svabhāv śānt nathī} — \textquotedblleft{}His temperament is not calm.\textquotedblright{}
\end{quote}

\begin{cousinweb}
\textbf{\gu{સ્વભાવ}} (\emph{svabhāv}) is Sanskrit \emph{svabhāva}, \textquotedblleft{}own being, nature\textquotedblright{}, built from \emph{sva}, \textquotedblleft{}own\textquotedblright{}, and \emph{bhāva}, \textquotedblleft{}being\textquotedblright{}, and taken over as a tatsama; EtymDB lists it as borrowed from Sanskrit (EtymDB).
\end{cousinweb}

\begin{grammarlens}[title={What you've built}]
One more word for everyday life at A2.
\end{grammarlens}

\subsection*{Your turn}
\begin{itemize}
\raggedright
  \item \emph{Say it:} \emph{svabhāv}
  \item \emph{Say it:} \emph{svabhāv}, once more
  \item \emph{Say it:} \emph{teno svabhāv śānt nathī} — His temperament is not calm.
  \item \emph{Recall:} say the Gujarati for to remember fondly, then the Gujarati for respect, then say \emph{svabhāv} again
\end{itemize}

\begin{tcolorbox}[breakable,colback=teal!4,colframe=teal!35!black,title={Before you move on}]
What does \gu{સ્વભાવ} mean? (\textquotedblleft{}Temperament, nature\textquotedblright{}.) Say it once more.
\end{tcolorbox}
\section[\texorpdfstring{\gu{દયાળુ} (dayāḷu)}{દયાળુ (dayāḷu)}]{\gu{દયાળુ} (dayāḷu) — kind, compassionate}

\label{lesson:GU-C270-dayalu}



\begin{quote}
Before the new one: say the Gujarati for respect, then the Gujarati for temperament.
\end{quote}

\subsection*{You'll want to know: \gu{દયાળુ}}
\textbf{\gu{દયાળુ}} — \emph{dayāḷu} — \textquotedblleft{}kind, compassionate\textquotedblright{}.

\begin{quote}
\textbf{\gu{શું} \gu{તમારા} \gu{સાહેબ} \gu{દયાળુ} \gu{છે}?} — \emph{śũ tamārā sāheb dayāḷu che?} — \textquotedblleft{}Is your boss kind?\textquotedblright{}
\end{quote}

\begin{cousinweb}
\textbf{\gu{દયાળુ}} (\emph{dayāḷu}) is built on \emph{dayā}, Sanskrit \emph{day\'{\={a}}}, \textquotedblleft{}sympathy\textquotedblright{} (CDIAL 6185), with the adjective suffix \emph{-āḷu}, \textquotedblleft{}full of\textquotedblright{}.
\end{cousinweb}

\begin{grammarlens}[title={What you've built}]
One more word for everyday life at A2.
\end{grammarlens}

\subsection*{Your turn}
\begin{itemize}
\raggedright
  \item \emph{Say it:} \emph{dayāḷu}
  \item \emph{Say it:} \emph{dayāḷu}, once more
  \item \emph{Say it:} \emph{śũ tamārā sāheb dayāḷu che?} — Is your boss kind?
  \item \emph{Recall:} say the Gujarati for respect, then the Gujarati for temperament, then say \emph{dayāḷu} again
\end{itemize}

\begin{tcolorbox}[breakable,colback=teal!4,colframe=teal!35!black,title={Before you move on}]
What does \gu{દયાળુ} mean? (\textquotedblleft{}Kind, compassionate\textquotedblright{}.) Say it once more.
\end{tcolorbox}
\section[\texorpdfstring{\gu{પ્રામાણિક} (prāmāṇik)}{પ્રામાણિક (prāmāṇik)}]{\gu{પ્રામાણિક} (prāmāṇik) — honest}

\label{lesson:GU-C270-pramanik}



\begin{quote}
Before the new one: say the Gujarati for temperament, then the Gujarati for kind.
\end{quote}

\subsection*{You'll want to know: \gu{પ્રામાણિક}}
\textbf{\gu{પ્રામાણિક}} — \emph{prāmāṇik} — \textquotedblleft{}honest\textquotedblright{}.

\begin{quote}
\textbf{\gu{તે} \gu{વેપારી} \gu{પ્રામાણિક} \gu{નથી}.} — \emph{te vepārī prāmāṇik nathī} — \textquotedblleft{}That trader is not honest.\textquotedblright{}
\end{quote}

\begin{cousinweb}
\textbf{\gu{પ્રામાણિક}} (\emph{prāmāṇik}) is Sanskrit \emph{prāmāṇika}, \textquotedblleft{}founded on authority\textquotedblright{}, built on \emph{pramāṇa}, \textquotedblleft{}measure, proof\textquotedblright{}, and taken over as a tatsama (CDIAL 8949).
\end{cousinweb}

\begin{grammarlens}[title={What you've built}]
One more word for everyday life at A2.
\end{grammarlens}

\subsection*{Your turn}
\begin{itemize}
\raggedright
  \item \emph{Say it:} \emph{prāmāṇik}
  \item \emph{Say it:} \emph{prāmāṇik}, once more
  \item \emph{Say it:} \emph{te vepārī prāmāṇik nathī} — That trader is not honest.
  \item \emph{Recall:} say the Gujarati for temperament, then the Gujarati for kind, then say \emph{prāmāṇik} again
\end{itemize}

\begin{tcolorbox}[breakable,colback=teal!4,colframe=teal!35!black,title={Before you move on}]
What does \gu{પ્રામાણિક} mean? (\textquotedblleft{}Honest\textquotedblright{}.) Say it once more.
\end{tcolorbox}
\section[\texorpdfstring{\gu{રમૂજી} (ramūjī)}{રમૂજી (ramūjī)}]{\gu{રમૂજી} (ramūjī) — funny, witty}

\label{lesson:GU-C270-ramuji}



\begin{quote}
Before the new one: say the Gujarati for kind, then the Gujarati for honest.
\end{quote}

\subsection*{You'll want to know: \gu{રમૂજી}}
\textbf{\gu{રમૂજી}} — \emph{ramūjī} — \textquotedblleft{}funny, witty\textquotedblright{}.

\begin{quote}
\textbf{\gu{મારા} \gu{કાકા} \gu{રમૂજી} \gu{છે}, \gu{પણ} \gu{મારા} \gu{પિતા} \gu{નથી}.} — \emph{mārā kākā ramūjī che, paṇ mārā pitā nathī} — \textquotedblleft{}My uncle is funny, but my father is not.\textquotedblright{}
\end{quote}

\begin{cousinweb}
\textbf{\gu{રમૂજી}} (\emph{ramūjī}) is built on \emph{ramūj}, \textquotedblleft{}jest\textquotedblright{}, with \emph{-ī}. \emph{Ramūj} is usually traced to Arabic \emph{rumūz}, plural of \emph{ramz}, \textquotedblleft{}a sign, a hint\textquotedblright{}, through Persian; its Arabic root is \emph{r-m-z} (Lane).
\end{cousinweb}

\begin{grammarlens}[title={What you've built}]
One more word for everyday life at A2.
\end{grammarlens}

\subsection*{Your turn}
\begin{itemize}
\raggedright
  \item \emph{Say it:} \emph{ramūjī}
  \item \emph{Say it:} \emph{ramūjī}, once more
  \item \emph{Say it:} \emph{mārā kākā ramūjī che, paṇ mārā pitā nathī} — My uncle is funny, but my father is not.
  \item \emph{Recall:} say the Gujarati for kind, then the Gujarati for honest, then say \emph{ramūjī} again
\end{itemize}

\begin{tcolorbox}[breakable,colback=teal!4,colframe=teal!35!black,title={Before you move on}]
What does \gu{રમૂજી} mean? (\textquotedblleft{}Funny, witty\textquotedblright{}.) Say it once more.
\end{tcolorbox}
\section[\texorpdfstring{\gu{ઘમંડી} (ghamanḍī)}{ઘમંડી (ghamanḍī)}]{\gu{ઘમંડી} (ghamanḍī) — arrogant, proud}

\label{lesson:GU-C270-ghamandi}



\begin{quote}
Before the new one: say the Gujarati for honest, then the Gujarati for funny.
\end{quote}

\subsection*{You'll want to know: \gu{ઘમંડી}}
\textbf{\gu{ઘમંડી}} — \emph{ghamanḍī} — \textquotedblleft{}arrogant, proud\textquotedblright{}.

\begin{quote}
\textbf{\gu{શું} \gu{તે} \gu{છોકરો} \gu{ઘમંડી} \gu{છે}?} — \emph{śũ te chokro ghamanḍī che?} — \textquotedblleft{}Is that boy arrogant?\textquotedblright{}
\end{quote}

\begin{cousinweb}
The origin of \textbf{\gu{ઘમંડી}} (\emph{ghamanḍī}) is not settled. It is built on \emph{ghamanḍ}, \textquotedblleft{}pride\textquotedblright{}, with \emph{-ī}, but CDIAL has no entry for \emph{ghamanḍ} and our other sources do not trace it.
\end{cousinweb}

\begin{grammarlens}[title={What you've built}]
One more word for everyday life at A2.
\end{grammarlens}

\subsection*{Your turn}
\begin{itemize}
\raggedright
  \item \emph{Say it:} \emph{ghamanḍī}
  \item \emph{Say it:} \emph{ghamanḍī}, once more
  \item \emph{Say it:} \emph{śũ te chokro ghamanḍī che?} — Is that boy arrogant?
  \item \emph{Recall:} say the Gujarati for honest, then the Gujarati for funny, then say \emph{ghamanḍī} again
\end{itemize}

\begin{tcolorbox}[breakable,colback=teal!4,colframe=teal!35!black,title={Before you move on}]
What does \gu{ઘમંડી} mean? (\textquotedblleft{}Arrogant, proud\textquotedblright{}.) Say it once more.
\end{tcolorbox}
